
Machine learning researchers publish thousands of datasets yearly, yet critical metadata fields remain systematically undocumented: HuggingFace datasets show 55-66\% presence for preprocessing code and data provenance fields, while UCI ML Repository datasets show 0-15\% presence for the same fields --- a 45-61 percentage point gap. A researcher attempting to reproduce a UCI dataset's preprocessing workflow finds only methodology text with no executable code, no dependency specifications, and no versioning metadata.

Prior work has documented documentation heterogeneity within platforms. Yang et al. (2024) analyzed HuggingFace dataset cards at scale (7,433 datasets), revealing marked heterogeneity in completion rates: Dataset Description and Structure sections show higher completion than Considerations sections, with completion correlating with dataset popularity. Strecker (2026) identified metadata conflicts across 8 geoscience and social science repositories, attributing incomplete DataCite metadata to implementation workflows and inter-standard schema differences. However, no study has quantitatively linked repository UX features to metadata completeness patterns across platforms.

This work introduces \textbf{friction score} (0-4 scale: automated field extraction + pre-filled templates + validation feedback + programmatic API access) as a quantifiable UX metric. Platforms with higher friction scores (HuggingFace friction=3: templates + validation + API) show 55-66\% optional field presence versus UCI (friction=0: manual web forms only) showing 0-15\%, with effect sizes of 45-61 percentage points. Required fields (license, version) show 75-95\% presence across platforms with weaker friction effects (12-20pp, Cramér's V 0.1-0.2), validating enforcement as a distinct mechanism.

The contributions of this work are:

\textbf{Cross-platform friction analysis:} Extension of Yang (2024)'s single-platform analysis to 3 ML repositories (OpenML, HuggingFace, UCI) at 10,000+ dataset scale, introducing friction score as explanatory variable.

\textbf{Mechanism distinction:} Demonstration that enforcement sets floor (UCI 74-75\%, HuggingFace/OpenML 89-95\%) while friction reduction influences optional field presence (45-61pp effect sizes). Gradient effects validated for 2 of 3 optional fields (data\_source\_url, collection\_date show UCI < OpenML < HuggingFace).

\textbf{Quantitative effect sizes:} Measurement of 45-61pp differences for optional fields and 12-20pp for required fields. Linear gradient for provenance fields (URLs, dates) suggests cumulative UX effects. Threshold for code fields suggests platform feature requirements create binary enablement.

The remainder of this paper is organized as follows: Section 2 reviews related work on dataset documentation and metadata standardization. Section 3 describes friction scoring methodology and large-scale extraction approach. Section 4 details experimental design. Section 5 presents results. Section 6 discusses mechanism interpretation and limitations. Section 7 concludes.
